\chapter{L'Observatoire : donner à voir}
\label{ch:observatoire}

\questionchapitre{Un entrepôt ne se consulte pas. Comment rendre la donnée consolidée
immédiatement lisible par ceux qui en ont l'usage, sans enfermer les méthodes développées dans
l'application qui les héberge ?}

\section{Une contrainte d'architecture posée d'emblée}

La plateforme comporte trois couches : une interface web en React~\cite{react}, un service
applicatif FastAPI~\cite{fastapi}, et une bibliothèque métier en Python. La contrainte structurante porte sur la troisième : \textbf{la
bibliothèque métier ne dépend ni du service applicatif ni de l'interface web}. Elle produit des
figures, donc dépend d'une bibliothèque de tracé, mais rien dans son code ne suppose le
navigateur ni le serveur.

Cette règle a une conséquence pratique déterminante pour un travail de recherche. Le code de
préparation des données, de calibration des modèles métier, d'entraînement, d'explicabilité et
de génération de contrefactuels s'exécute indifféremment depuis le service web, depuis un script
ou depuis un carnet de calcul. Les méthodes développées ne sont donc pas prisonnières de
l'application : si l'interface était abandonnée, elles resteraient utilisables.

\begin{figure}[H]
\centering
\begin{tikzpicture}[
  font=\sffamily,
  bande/.style={draw=black!70, fill=white, minimum width=10.6cm, minimum height=0.85cm,
                align=center, font=\sffamily\footnotesize},
  coeur/.style={draw=black, very thick, fill=black!3, minimum width=10.6cm,
                minimum height=1.9cm, align=center, font=\sffamily\footnotesize},
  mod/.style={draw=black!50, fill=white, minimum width=2.35cm, minimum height=0.7cm,
              align=center, font=\sffamily\scriptsize},
  res/.style={draw=black!55, fill=white, minimum width=3.3cm, minimum height=0.8cm,
              align=center, font=\sffamily\scriptsize},
  fl/.style={-{Stealth[length=2.2mm]}, draw=black!70, thick},
  note/.style={font=\sffamily\scriptsize\itshape, color=black!60},
]

\node[bande] (spa) {\textbf{Interface web} \\[1pt]
                    \scriptsize cartes, graphiques, suivi d'entraînement};
\node[bande, below=0.75cm of spa] (api) {\textbf{Service applicatif} \\[1pt]
                    \scriptsize API et flux d'événements};
\node[coeur, below=0.75cm of api] (core) {};
\node[anchor=north, font=\sffamily\footnotesize\bfseries, yshift=-3pt] at (core.north)
     {Bibliothèque métier};
\node[anchor=north, font=\sffamily\scriptsize, yshift=-17pt] at (core.north)
     {aucune dépendance d'interface};

\node[mod] (m1) at ($(core.center)+(-3.75,-0.42)$) {Préparation\\des données};
\node[mod] (m2) at ($(core.center)+(-1.25,-0.42)$) {Modèles\\métier};
\node[mod] (m3) at ($(core.center)+(1.25,-0.42)$)  {Prévision\\profonde};
\node[mod] (m4) at ($(core.center)+(3.75,-0.42)$)  {Explicabilité\\contrefactuels};

\node[res, below=0.95cm of core, xshift=-3.65cm] (cache) {Cache de réponses};
\node[res, below=0.95cm of core, draw=black, very thick] (entrepot)
     {Entrepôt \\[1pt] \scriptsize lecture SQL};
\node[res, below=0.95cm of core, xshift=3.65cm] (registre) {Registre d'expériences};

\draw[fl] (spa) -- (api);
\draw[fl] (api) -- (core);
\draw[fl] (core.south) -- (entrepot.north);
\draw[fl] (core.south west) ++(1.1,0) -- ++(0,-0.3) -| (cache.north);
\draw[fl] (core.south east) ++(-1.1,0) -- ++(0,-0.3) -| (registre.north);

\node[note, below=0.14cm of entrepot] {dépôt de l'entrepôt, chapitre~\ref{ch:entrepot}};

\end{tikzpicture}
\caption{Architecture de la plateforme. La bibliothèque métier est appelable depuis le service
web, un script ou un carnet de calcul, sans modification. C'est cette propriété qui rend les
développements réutilisables au-delà de l'application.}
\label{fig:archi-plateforme}
\end{figure}

\section{Le besoin : explorer avant de modéliser}

Un entrepôt bien construit reste inerte tant qu'il n'est pas regardé. Deux publics le
demandaient, pour des raisons différentes.

Les \textbf{hydrogéologues} ont besoin de situer une station dans son contexte : quelle nappe
observe-t-elle, comment se comporte-t-elle par rapport à ses voisines, où en est-elle dans son
historique, et que faisait le climat pendant ce temps. Ce travail de mise en regard, fait
manuellement, mobilise plusieurs sources et plusieurs outils.

Les \textbf{modélisateurs} ont besoin de choisir leurs cas d'étude en connaissance de cause :
une station dont la chronique est trop courte, trop lacunaire ou visiblement perturbée par un
pompage ne fera pas un bon cas d'apprentissage. Repérer ces situations avant d'entraîner un
modèle évite d'attribuer à la méthode ce qui relève de la donnée.

L'Observatoire répond aux deux.

\section{Organisation : deux compartiments}

L'interface s'organise en deux volets correspondant aux deux compartiments du cycle observé.

Le premier volet, nappes et rivières, donne l'exploration spatiale des stations piézométriques
et hydrométriques sur fond cartographique, avec superposition des entités aquifères du référentiel
BDLISA et d'un découpage en secteurs hydrogéologiques. Les indicateurs de situation y sont
portés directement par les stations, ce qui donne une lecture immédiate de l'état du territoire.
La comparaison entre stations est possible, avec conservation de la sélection d'une vue à
l'autre. La figure~\ref{fig:obs-carte-nationale} en donne la vue d'ensemble, la
figure~\ref{fig:obs-station} la lecture d'un ouvrage particulier.

\begin{figure}[htbp]
\centering
\vue[0.97]{carte-nationale}
\caption{Le volet nappes et rivières à l'échelle du pays. Chaque pastille agrège les stations
d'un groupe et porte la couleur de leur classe de situation ; le curseur du bas rejoue la période
2000--2026 saison par saison. Le compteur donne l'assiette réelle du classement, \num{2610}
piézomètres classés sur \num{22407} recensés et \num{2170} stations hydrométriques sur
\num{6253} : les stations non classées le sont faute d'historique suffisant pour standardiser, et
l'écart est visible plutôt que masqué.}
\label{fig:obs-carte-nationale}
\end{figure}

\begin{figure}[htbp]
\centering
\vue[0.97]{carte-regionale}
\caption{La même carte à l'échelle régionale, une station sélectionnée. Le panneau de gauche
porte l'indicateur de situation, la position de la mesure du mois sur les sept classes de
référence, la profondeur d'historique, et les stations de la même entité hydrogéologique BDLISA
avec leur propre classe. C'est cette mise en regard, faite manuellement sur plusieurs sources,
que le volet remplace par une sélection. \lecas}
\label{fig:obs-station}
\end{figure}

Le second, le climat, donne l'exploration de la grille nationale de \ang{0.1} selon trois
vues : une vue de
situation cartographique à l'échelle du pays, une vue par point ou par zone donnant l'historique
d'une maille depuis 1950 (pluie comparée à la normale, indices sur plusieurs fenêtres, table des
épisodes de sécheresse, export tabulaire), et une vue de comparaison inter-annuelle.

Les deux volets communiquent : depuis la carte des nappes, un lien contextuel ouvre le climat de
la maille correspondante ; depuis la fiche d'une station, on accède à son contexte climatique
local. La figure~\ref{fig:obs-climat} montre ce croisement sur une station, la
figure~\ref{fig:obs-comparaison} la comparaison inter-annuelle.

\begin{figure}[htbp]
\centering
\vue[0.97]{contexte-spi}
\caption{Le contexte climatique d'une station, sur la même page que sa chronique : indice
standardisé de précipitation sur trois mois depuis dix ans, coloré selon les sept classes, et
cumuls glissants comparés à la normale 1991--2020. Les valeurs affichées sont celles produites et
validées au chapitre~\ref{ch:indicateurs}, non un recalcul.}
\label{fig:obs-climat}
\end{figure}

\begin{figure}[htbp]
\centering
\vue[0.97]{comparaison-annees}
\caption{La comparaison inter-annuelle : plusieurs années superposées mois par mois sur le même
ouvrage, à sélectionner dans la liste des années disponibles. La lecture recherchée est
qualitative, situer l'année courante parmi les années remarquables, et elle ne suppose aucun
seuillage inventé.}
\label{fig:obs-comparaison}
\end{figure}

\section{Une règle de conception qui a structuré l'interface}

Les couches cartographiques disponibles obéissent à une règle explicite, adoptée après
discussion et maintenue depuis :

\begin{quote}
\emph{Soit un indicateur réellement standardisé, soit une valeur réelle affichée comme telle.
Jamais un intermédiaire inventé.}
\end{quote}

Concrètement, les indices standardisés (précipitation, température, bilan hydrique) sont
cartographiés sur l'échelle à sept classes, et les grandeurs journalières brutes (températures
extrêmes, pluie du jour) le sont sur une échelle absolue. En revanche, les cumuls mensuels
absolus de pluie, de température ou d'évapotranspiration ne sont \textbf{pas} proposés comme
couches cartographiques : ils apparaissent comme chiffres exacts dans le panneau de détail d'une
maille.

La raison est méthodologique. Cartographier une grandeur absolue sur une palette de couleurs
suppose un seuillage, donc une norme implicite. Or cette norme, n'étant pas une référence
statistique établie, serait une invention de l'outil. L'utilisateur lirait une carte de
sécheresse là où il n'y a qu'un choix de bornes. Refuser cette couche coûte une fonctionnalité
apparente et évite une erreur d'interprétation systématique.

\section{Trois choix techniques}

L'Observatoire recalcule le moins possible : ses points d'accès lisent les tables analytiques de l'entrepôt plutôt que de refaire le calcul. La carte, la
situation courante et les exports lisent ainsi les indicateurs déjà produits, de sorte que
\textbf{la valeur affichée est exactement celle qui a été validée} au
chapitre~\ref{ch:indicateurs}. Trois points d'accès font exception : les deux qui servent la
série historique d'une station, et celui qui calcule l'antériorité portée par la carte de
situation. Ils lisent la grille de référence dans l'entrepôt, mais appliquent la standardisation
à la requête. Le calcul reste le même, et un test croisé garde la cohérence des classes entre
les deux dépôts.

Les réponses sont mises en cache, ensuite. Les vues cartographiques nationales portent sur des
volumes importants et changent peu ; leur mise en cache pour vingt-quatre heures réduit la
charge sans affecter la fraîcheur perçue. Les couches qui bougent plus vite ont des durées plus
courtes, de quelques heures pour le fond de carte des stations à une heure pour les alertes. Cette optimisation a un corollaire opérationnel
documenté : tout déploiement qui modifie la forme ou la valeur des réponses doit purger le
cache, sous peine de servir des valeurs périmées pendant une journée.

Le point d'entrée, enfin, est unique. Le service qui sert l'interface assure également le relais
vers l'API, les en-têtes de sécurité et la limitation de débit. Regrouper ces fonctions dans un
seul composant vaut moins pour la simplicité que pour ce que cela écarte : une protection
configurée dans un composant qui ne tourne pas est une protection absente, avec en plus la
croyance qu'elle existe.

\section{Environnements et déploiement}

La séparation des environnements décrite au chapitre~\ref{ch:entrepot} prend ici tout son sens :
elle autorise à expérimenter sur une chaîne en évolution sans compromettre l'instance qui sert
les utilisateurs. La production est répartie : l'interface est déployée sur l'infrastructure
mutualisée de l'établissement, le calcul et les bases restent sur le serveur équipé d'un
processeur graphique. Une seule exigence réseau conditionne l'ensemble, l'ouverture de la route
entre les deux.

L'Observatoire dépend entièrement de l'entrepôt, et ne se dégrade pas gracieusement en son
absence : il répond en erreur plutôt que d'afficher un état partiel. Le choix est assumé pour un
outil scientifique, où une valeur absente vaut mieux qu'une valeur douteuse. Il serait à revoir
pour une diffusion large.
